\section{Introduction}
\label{sec:intro}
% 1-2 страницы
% why LLMs need reliable forgetting (privacy, dynamic world knowledge)
% LLM Unlearning
% LLM knolage saving: Popular vs rare facts present different challenges (surface‑level leakage vs data sparsity)
% Our contributions


Large Language Models (LLMs) are central to modern AI, yet their training on massive web datasets often embeds private, copyrighted, or erroneous information \cite{wei2024evaluating, jang2023knowledge, NEURIPS2024_b1f78dfc}.
Removing such material after a model has been trained, known as \textit{machine unlearning}, is therefore both a legal necessity and a technical challenge \cite{ren2025keeping}.
The goal is to excise a targeted subset of training data so thoroughly that the model behaves as if those examples were never seen, while preserving its overall competence and avoiding the prohibitive cost of full retraining \cite{geng2025comprehensive, jang2023knowledge, zagardo2024more}.

Despite rapid progress, current unlearning algorithms share an assumption: \textbf{all facts are equally forgettable} \cite{geng2025comprehensive, hu2024unlearning}.  
In practice, knowledge appears in training data with wildly different frequencies.  
A canonical fact such as \enquote{Paris is the capital of France} surfaces thousands of times across diverse sources, whereas an obscure biographical detail may occur only once.  

Intuitively, the former presents a significantly greater unlearning challenge due to its extensive embedding; however, current state-of-the-art unlearning methodologies often apply a uniform \enquote{unlearning force} irrespective of knowledge prevalence or entrenchment \cite{ren2025keeping, NEURIPS2024_16e18fa3}.
This mismatch can yield two failure modes:
\begin{itemize}
    \item \textbf{Ineffective forgetting} – the disfavoured fact remains recoverable;
    \item \textbf{Catastrophic forgetting} – collateral erasure of unrelated knowledge that degrades the model’s utility.
\end{itemize}

Examples of ineffective and effective unlearning of a popular fact are illustrated in Figure~\ref{fig:mu_modes}. 
% Thus, \noindent\textbf{knowledge popularity} emerges as a critical, but still unexplored, dimension in assessing and designing machine‑unlearning techniques \cite{chang2024large}.

Although the two common failure modes, such as ineffective forgetting and catastrophic forgetting, stem from applying a uniform unlearning force despite the large differences in how frequently facts appear in pretraining, the relationship between fact popularity and these failure modes has not been studied.

% Studies of fact popularity effects have focused on QA and relation extraction, showing that head facts are retained and retrieved more reliably than long-tail ones in language models, creating systematic gaps \cite{petroni2019language, kandpal2023large, ghosal2024understanding, sun2024head, henning2023survey}, yet in MU the interaction between fact popularity and common failure modes (e.g., ineffective forgetting vs. catastrophic forgetting from applying a uniform unlearning force) remains unexplored.


%\paragraph{Our Approach.}  
% We present the first systematic study of how \emph{knowledge popularity} interacts with \emph{model scale} in machine unlearning.  
% We base our benchmark, \textbf{UNLamb}, on the original PopQA \cite{mallen2023not} corpus, which already provides per‑fact popularity scores (\texttt{s\_pop}).  
% We re‑structure PopQA into unlearning‑style splits—matched \emph{rare} and \emph{popular} forget sets at multiple ratios—keeping the remaining data for retain‑ and validation‑splits.  
% This adaptation makes PopQA directly usable for controlled unlearning experiments.
In this paper, we start from PopQA, an open‑domain, entity‑centric QA dataset derived from Wikidata triples (14.3k triplets); each item corresponds to a subject–relation pair and is annotated with a popularity score (\texttt{s\_pop}) computed from Wikipedia page views.  
Leveraging these annotations, the novel benchmark, \textbf{UNLamb} \textit{-- UNlearning LAnguage Models Benchmark}, is introduced, in which we stratify questions by popularity, forming unlearning-style splits—matched \emph{rare}/\emph{popular} forget sets, and keep the remainder for retain/validation sets.
Because Wikidata is included in the pretraining corpora of many LLMs, our benchmark approximates real-world knowledge removal more closely than synthetic alternatives. Using UNLamb, we evaluate representative unlearning algorithms on LLaMA-3 family models from 1B to 8B parameters.  
Our key finding is that, as model size grows, attempts to erase \emph{popular}, deeply embedded facts raise the risk of dataset‑wide catastrophic forgetting; in contrast, removing similarly sized sets of \emph{rare} facts is slightly more likely to succeed with less collateral damage.  
Importantly, even rare‑fact removal is \textbf{not} immune to failure—ineffective forgetting or collateral loss still occurs, albeit less often, so hyper‑parameters must be tuned more carefully and the popularity range in the forget set must be taken into account.  
While the popularity effect is muted in smaller models, it becomes increasingly pronounced at larger scales, underscoring the need for popularity‑aware, scale‑adaptive unlearning strategies.

\paragraph{Our main contributions and findings:}
\begin{enumerate}
    \item Propose knowledge popularity as a critical feature for unlearning algorithms and articulate the mechanisms through which it impacts unlearning efficacy.
    \item Publish \textbf{UNLamb}, a novel, popularity-stratified QA benchmark for a more precise evaluation of unlearning methods, complemented by a public repository to ensure full reproducibility.
    \item Conduct large-scale experiments on the LLaMA-3 family and find that popular knowledge is markedly more difficult to unlearn, an effect that is exacerbated as model size increases.
    % \item We propose \emph{knowledge popularity} as a important feature for applying machine unlearning algorithms and explain why it matters for effective unlearning.
    % \item We \emph{publish} \textsc{UNLamb}, a popularity‑stratified QA benchmark for more precise evaluation of unlearning methods, and provide an open Git repository for full reproducibility.
    % \item We \textit{conduct broad experiments} on three sizes of the LLaMA 3 family using state-of-the-art unlearning techniques and find that popular facts are markedly harder to erase, with the effect intensifying as model size grows.
\end{enumerate}


% The UNLamb benchmark is available on the Hugging Face Hub:
% \url{https://huggingface.co/datasets/SwetieePawsss/UNLamb}.

% Repository with experiments on UNLamb is hosted on GitHub with anonymization: \url{https://anonymous.4open.science/r/UNLamb_unlearning-B7E4/README.md}.